% F085｜真实数学图；代表函数不构成抽象解的唯一性证明。
\begin{center}
\begin{tikzpicture}
\begin{axis}[width=12.4cm,height=5.1cm,axis lines=middle,
 xmin=-2.5,xmax=2.6,ymin=-0.5,ymax=4.8,
 xtick={-2,-1,0,1,2},xticklabels=\empty,ytick={0,1,2,4},xlabel={$t$},ylabel={$f(t)$},
 title={指数型的两个分支：非零代表与零函数},title style={font=\small,yshift=3pt},
 tick label style={font=\small},label style={font=\small},
 ylabel style={at={(axis description cs:0,1)},anchor=south east,xshift=-18pt},
 clip=false,grid=major,grid style={black!8},enlargelimits=false,
 every axis plot/.append style={thick},samples=161]
\addplot[CExplanation,very thick,domain=-2:2.1] {2^x};
\addplot[orange!80!black,dashed,domain=-2:2.1] {0};\addplot[only marks,mark=*,mark size=2pt,orange!80!black] coordinates {(0,1) (0,0)};

\node[font=\small,anchor=north,yshift=-3pt] at (axis cs:-2,-0.5) {$-2$};
\node[font=\small,anchor=north,yshift=-3pt] at (axis cs:-1,-0.5) {$-1$};
\node[font=\small,anchor=north,yshift=-3pt] at (axis cs:1,-0.5) {$1$};
\node[font=\small,anchor=north,yshift=-3pt] at (axis cs:2,-0.5) {$2$};

\end{axis}
\end{tikzpicture}
\par\small 蓝实线 $2^t$ 与橙虚线 $0$ 都满足 $f(x+y)=f(x)f(y)$，在 $0$ 处的值分别为 $1,0$；先排除零函数才能约去非零值。
\end{center}
